% Module 10: QC Hardware Platforms and Cloud Access
% Estimated slides: 25 | Duration: ~2 hours
% Key topics: IBM Quantum platform, Qiskit architecture, transpilation, noise models,
%             Google Cirq, Amazon Braket, Azure Quantum, platform comparison
% Labs: End-to-end IBM Quantum submission; noise characterization with Qiskit Aer

%%%%%%%%%%%%%%%%%%%%%%%%%%%%%%%%%%%%%%%%%%%%%%%%%%%%%%%%%%
\begin{frame}[fragile]\frametitle{}
\begin{center}
{\Large Module 10: QC Hardware Platforms}
\end{center}
\end{frame}

%%%%%%%%%%%%%%%%%%%%%%%%%%%%%%%%%%%%%%%%%%%%%%%%%%%%%%%%%%%
\begin{frame}[fragile]\frametitle{Cloud Quantum Computing: Access Today}
\begin{itemize}
\item You do not need a dilution refrigerator to use a quantum computer
\item Cloud-based quantum computing: access real QPUs over the internet
\item Major platforms: IBM Quantum, Google Quantum AI, Amazon Braket, Azure Quantum
\item Free access: IBM Quantum (limited); paid plans for priority and more QPUs
\item Workflow: write circuit in Python $\to$ submit to cloud $\to$ receive results
\item Waiting time: seconds to minutes depending on queue length
\end{itemize}
% Suggested diagram: Cloud access workflow from laptop to quantum hardware and back
\end{frame}

%%%%%%%%%%%%%%%%%%%%%%%%%%%%%%%%%%%%%%%%%%%%%%%%%%%%%%%%%%%
\begin{frame}[fragile]\frametitle{IBM Quantum Platform: Overview}
\begin{itemize}
\item IBM has offered cloud quantum access since 2016 (5 qubits initially)
\item Current offering: 20+ systems globally; largest open-access: 127 qubit Eagle
\item IBM Quantum Network: partners (research, industry) get premium access
\item Composer: visual circuit builder; Quantum Lab: Jupyter notebook environment
\item Qiskit Runtime: optimized execution service; reduces round-trip overhead
\item IBM roadmap: 1000+ qubits by 2023 (Condor/Heron), error correction by 2029
\end{itemize}
% Suggested diagram: IBM Quantum website screenshot layout (anonymized); chip photo
\end{frame}

%%%%%%%%%%%%%%%%%%%%%%%%%%%%%%%%%%%%%%%%%%%%%%%%%%%%%%%%%%%
\begin{frame}[fragile]\frametitle{Qiskit Architecture}
\begin{itemize}
\item Qiskit: IBM's open-source quantum SDK; most widely used quantum framework
\item \textbf{qiskit}: core circuit library and transpiler
\item \textbf{qiskit-aer}: high-performance classical simulators
\item \textbf{qiskit-ibm-runtime}: interface to IBM Quantum hardware
\item \textbf{qiskit-nature}: quantum chemistry and physics
\item \textbf{qiskit-machine-learning}: quantum ML algorithms
\item \textbf{qiskit-finance}: quantum finance (options pricing, portfolio optimization)
\item \textbf{qiskit-optimization}: combinatorial optimization (QAOA)
\end{itemize}
% Suggested diagram: Qiskit ecosystem package hierarchy
\end{frame}

%%%%%%%%%%%%%%%%%%%%%%%%%%%%%%%%%%%%%%%%%%%%%%%%%%%%%%%%%%%
\begin{frame}[fragile]\frametitle{Transpilation: From Your Circuit to Hardware}
\begin{itemize}
\item Your circuit uses arbitrary gates; hardware only supports a native gate set
\item Transpilation: converts your circuit to hardware-native gates and connectivity
\item Steps: unroll to basis gates $\to$ map to physical qubits $\to$ route/SWAP $\to$ optimize
\item Optimization levels: 0 (no optimization) to 3 (heavy optimization)
\item SWAP overhead: if qubits that interact are not connected, insert SWAP gates
\item Check circuit depth before and after transpilation --- can increase significantly
\end{itemize}
\begin{verbatim}
from qiskit import transpile
transpiled = transpile(qc, backend=backend,
                       optimization_level=3)
print(transpiled.depth())
\end{verbatim}
% Suggested diagram: Coupling map of a real device; circuit before vs after transpilation
\end{frame}

%%%%%%%%%%%%%%%%%%%%%%%%%%%%%%%%%%%%%%%%%%%%%%%%%%%%%%%%%%%
\begin{frame}[fragile]\frametitle{Qiskit Aer: Simulators}
\begin{itemize}
\item Aer provides fast classical simulators for testing circuits before hardware runs
\item \textbf{StatevectorSimulator}: exact simulation; tracks full $2^n$ state vector
\item \textbf{QasmSimulator}: samples measurement outcomes (mimics real hardware)
\item \textbf{AerSimulator}: unified simulator with noise model support
\item \textbf{NoiseModel}: add realistic noise (depolarizing, amplitude damping, readout error)
\item Noise models can be loaded from real device calibration data
\end{itemize}
\begin{verbatim}
from qiskit_aer.noise import NoiseModel
from qiskit_ibm_runtime import QiskitRuntimeService
noise_model = NoiseModel.from_backend(backend)
\end{verbatim}
\end{frame}

%%%%%%%%%%%%%%%%%%%%%%%%%%%%%%%%%%%%%%%%%%%%%%%%%%%%%%%%%%%
\begin{frame}[fragile]\frametitle{Lab: Running on IBM Quantum Hardware}
\begin{itemize}
\item Authenticate with IBM Quantum: \texttt{QiskitRuntimeService.save\_account(token=...)}
\item Choose a backend with fewest current jobs in queue
\item Transpile your Bell state circuit for the chosen backend
\item Submit as a Sampler job; wait for results
\item Compare hardware results vs ideal simulator results
\item Observe: hardware gives imperfect correlations due to noise
\end{itemize}
\begin{verbatim}
from qiskit_ibm_runtime import QiskitRuntimeService, SamplerV2
service = QiskitRuntimeService()
backend = service.least_busy(operational=True, simulator=False)
\end{verbatim}
\end{frame}

%%%%%%%%%%%%%%%%%%%%%%%%%%%%%%%%%%%%%%%%%%%%%%%%%%%%%%%%%%%
\begin{frame}[fragile]\frametitle{Google Quantum AI and Cirq}
\begin{itemize}
\item Google has developed Sycamore (53 qubits) and Willow (105 qubits) processors
\item Cirq: Google's open-source quantum framework (Python)
\item Access: Google Cloud Quantum AI; limited public access to real hardware
\item Focus: quantum supremacy demonstrations; surface code error correction research
\item Willow chip (2024): error rates below surface code threshold; exponential error reduction
\item Cirq designed for: quantum variational algorithms, circuit optimization
\end{itemize}
% Suggested diagram: Cirq circuit syntax example; Google chip photo
\end{frame}

%%%%%%%%%%%%%%%%%%%%%%%%%%%%%%%%%%%%%%%%%%%%%%%%%%%%%%%%%%%
\begin{frame}[fragile]\frametitle{Amazon Braket}
\begin{itemize}
\item AWS service: access multiple quantum hardware providers through one API
\item Supported hardware: IonQ (trapped ion), Rigetti (superconducting), OQC, Xanadu (photonic)
\item Also offers managed simulators (local, managed cloud)
\item PennyLane integration: quantum ML with hardware backends
\item Pay-per-task pricing: convenient for exploration without large commitment
\item SDK: \texttt{amazon-braket-sdk} (Python); familiar AWS workflow
\end{itemize}
% Suggested diagram: Braket multi-backend architecture showing different providers
\end{frame}

%%%%%%%%%%%%%%%%%%%%%%%%%%%%%%%%%%%%%%%%%%%%%%%%%%%%%%%%%%%
\begin{frame}[fragile]\frametitle{Microsoft Azure Quantum}
\begin{itemize}
\item Azure Quantum: cloud service integrating quantum hardware + optimization tools
\item Hardware partners: IonQ, Quantinuum (H-Series trapped ion), Pasqal
\item Q\# language: Microsoft's domain-specific language for quantum algorithms
\item QDK (Quantum Development Kit): Q\# + Python integration; resource estimator
\item Microsoft focus: topological qubits (long-term), Azure resource estimation
\item Unique tool: Azure Quantum Resource Estimator -- plan qubit requirements before hardware
\end{itemize}
% Suggested diagram: Azure Quantum architecture with hardware, simulator, and tools layers
\end{frame}

%%%%%%%%%%%%%%%%%%%%%%%%%%%%%%%%%%%%%%%%%%%%%%%%%%%%%%%%%%%
\begin{frame}[fragile]\frametitle{Platform Comparison for Practitioners}
\begin{itemize}
\item Best for learning: IBM Quantum (free access, best documentation, Qiskit ecosystem)
\item Best fidelity today: Quantinuum H2 (trapped ion, $> 99.9\%$ 2-qubit gates)
\item Most hardware variety: Amazon Braket (multiple vendors in one API)
\item Best for enterprise: Azure Quantum (enterprise support, resource estimation)
\item PennyLane (Xanadu): hardware-agnostic; best for quantum ML research
\item Bottom line: start with Qiskit/IBM Quantum; branch out as needed
\end{itemize}
% Suggested diagram: Side-by-side comparison table: Platform | Backend | SDK | Free tier | Best for
\end{frame}

%%%%%%%%%%%%%%%%%%%%%%%%%%%%%%%%%%%%%%%%%%%%%%%%%%%%%%%%%%%
\begin{frame}[fragile]\frametitle{Module 10 Summary}
\begin{itemize}
\item Cloud quantum access: IBM Quantum, Google AI, Braket, Azure --- no hardware needed
\item Qiskit ecosystem: core, Aer (simulators), Runtime (hardware), Nature/ML/Finance
\item Transpilation: converts your circuit to hardware-native gates and qubit layout
\item Simulation before hardware: always test on Aer first (faster, cheaper, more informative)
\item IBM Quantum free tier: best starting point for hands-on quantum programming
\item Noise models: load from real device data to simulate realistic hardware behavior
\item Next: Qiskit Programming Lab (hands-on full pipeline)
\end{itemize}
\end{frame}
